\clearpage\thispagestyle{empty}
\pdfbookmark[0]{Quellen zum Brief über den zweiten Kometen von 1805}{cometletter18060520-sources}
\begin{center}\textbf{Quellen zum Brief über den zweiten Kometen von 1805}\end{center}
\noindent Diese KI-Übertragung in heutiges Deutsch folgt dem Brief Carl Friedrich Gauß’ vom 20. Mai 1806 an Franz Xaver von Zach in \href{https://gdz.sub.uni-goettingen.de/id/PPN236007467}{Gauß’ \textit{Werken}, Band~VI} (Göttingen, 1874), S.~270--275. Der Brief erschien in der \textit{Monatlichen Correspondenz}, Band~XIV, Juli 1806, Beitrag~VIII, Teil~2, \href{https://www.digitale-sammlungen.de/en/view/bsb10538606?page=85}{S.~75}--\href{https://www.digitale-sammlungen.de/en/view/bsb10538606?page=96}{86}. Die Zeitschriftendigitalisate stammen von der Bayerischen Staatsbibliothek. Die Zeitspalte und die eckige Klammer zur Ortszeit folgen der Werkausgabe, in der auf S.~273 \textit{dasss} steht.
\par\medskip
Gauß verweist auf Jérôme Lalandes \textit{Histoire céleste française} (Paris, 1801), \href{https://archive.org/details/histoireclestef00lalagoog/page/n47/mode/1up}{S.~20}, \href{https://archive.org/details/histoireclestef00lalagoog/page/n64/mode/1up}{S.~37}, \href{https://archive.org/details/histoireclestef00lalagoog/page/n227/mode/1up}{S.~200} und \href{https://archive.org/details/histoireclestef00lalagoog/page/n516/mode/1up}{S.~477}, sowie auf die \textit{Connaissance des tems} für \href{https://cdt.imcce.fr/items/show/152}{das Jahr~XI} und \href{https://cdt.imcce.fr/items/show/154}{das Jahr~XIII}. Übertragung und Satz wurden mit KI erstellt.
\clearpage
